% ---------------------------------------------------------------------------
\subsection{Functional-Form Identifiability: Is GRPO Reward Growth Exponential?}
\label{sec:scaling-law-iter37}
% ---------------------------------------------------------------------------

\paragraph{Motivation.}
Iter~17--33 fit $R(t) = R_{\max}\,(1 - e^{-\lambda t})$ to the twelve-anchor
frontier set as the default saturation form, and showed that the
three-phase hypothesis (slow start, rapid improvement, plateau) is
\emph{partially falsified} on that set (P2 fails: 5/12 anchors are drift
or collapse; P3 borderline; P4 decisively falsified; only P5 -- Nemotron
uniqueness -- sustained, see \S\ref{sec:scaling-law-iter33}). None of
the prior iters, however, tested whether the exponential
\emph{itself} is the right functional form. If the literature's
exponential is mis-specified, the saturation ceiling estimates
$R_{\max}, \lambda, t_{80}$ propagated through the paper are biased.

\paragraph{Method.}
Iter~37 fits five candidate reward-curve forms to each anchor via
nonlinear least squares and ranks them by AIC and BIC with Akaike
weights:
\begin{description}
\item[(A)] Exponential saturation: $R_{\max}\,(1 - e^{-\lambda t})$
\item[(B)] Michaelis-Menten: $R_{\max}\,t/(t_{1/2} + t)$
\item[(C)] Hill ($n=2$): $R_{\max}\,t^2/(K^2 + t^2)$
\item[(D)] Power-law approach: $R_{\max}\,(1 - (1 + t)^{-\alpha})$
\item[(E)] Linear (null/baseline): $a + b\,t$
\end{description}
We run the battery three ways:
\begin{itemize}
\item \textbf{Synthetic (iter37 main, 12 anchors):} the same per-step
  trace reconstruction used in iter33, so the result is directly
  comparable to the iter33 phase battery.
\item \textbf{Dynamic-only (iter37b, 10 anchors):} filter to
  $\mathrm{var}(\mathrm{reward}) > 0.005$, removing the two flat
  anchors (Qwen3-30B-MoE-Inst, Qwen3-235B-MoE) that have no curvature
  signal by construction.
\item \textbf{Raw per-step cross-check (iter37c, 10 runs):} fit the
  same five forms to the 40-step per-step reward trajectories from
  the same-stack PPO/GRPO runs in
  \texttt{experiments/results/samestack\_ppo\_grpo.json}
  (5 GRPO + 5 PPO seeds, Qwen/Qwen2.5-0.5B).
\end{itemize}
We further run a parametric bootstrap (B=500, $\sigma=0.10$ added
reward noise) to estimate the win-rate of each form under sampling
variability, and a 2000-replicate column-permutation test of the
hypothesis ``$w_{\mathrm{sat}} = $ uniform(1/5)''.

\paragraph{Results: synthetic 12-anchor set (iter37).}
On the 12-anchor frontier set reconstructed from summary statistics,
the linear baseline wins AIC on 9/12 anchors; the exponential
saturation form (A) wins on 1/12 (Qwen3-30B-MoE). The mean Akaike
weight for the exponential is $\bar{w}_{\mathrm{sat}} = 0.044$, well
below the uniform $0.2$ prior. The bootstrap
($B{=}200$, Gaussian noise $\sigma{=}0.10$) gives a saturation
win-share of $0.096$ vs.\ a linear win-share of $0.732$.

\paragraph{Results: dynamic-only 10-anchor set (iter37b).}
The two flat anchors (Qwen3-30B-MoE-Inst, Qwen3-235B-MoE) are removed;
the remaining 10 anchors have measurable $\mathrm{var}>0.005$. Even
on this dynamic subset, linear wins 7/10, saturation 1/10, Hill 1/10,
Michaelis-Menten 1/10, power-law 0/10. Mean Akaike weights:
$\bar{w}_{\mathrm{linear}}=0.66$, $\bar{w}_{\mathrm{hill}}=0.13$,
$\bar{w}_{\mathrm{MM}}=0.10$, $\bar{w}_{\mathrm{sat}}=0.05$,
$\bar{w}_{\mathrm{power}}=0.06$. The bootstrap win-share is
$\mathrm{linear}=0.72$, $\mathrm{hill}=0.13$, $\mathrm{MM}=0.04$,
$\mathrm{sat}=0.09$. The 2000-replicate column-permutation test gives
$p=1.0$ for the hypothesis ``the saturation column carries more
weight than the average permuted column'' -- i.e.\ the data
\emph{rejects} the exponential as uniquely informative.

\paragraph{Results: raw per-step cross-check (iter37c).}
The synthetic-trace result is potentially an artefact: the
synthesised trace is too smooth to expose the curvature structure
of real per-step data. We therefore fit the same five forms to the
\emph{real} 40-step per-step reward trajectories from
\texttt{samestack\_ppo\_grpo.json} (Qwen2.5-0.5B, 5 seeds each of
GRPO and PPO). The result inverts: linear wins 0/5 on GRPO; the
Hill n=2 form wins 3/5, exponential wins 2/5, MM wins 0/5. On PPO,
Hill wins 2/5, with sat / linear / MM splitting the rest. Mean
Akaike weights on GRPO: $\bar{w}_{\mathrm{hill}}=0.62$,
$\bar{w}_{\mathrm{sat}}=0.38$, all other forms $\le 0.001$. On PPO,
$\bar{w}_{\mathrm{hill}}=0.40$, $\bar{w}_{\mathrm{MM}}=0.25$,
$\bar{w}_{\mathrm{sat}}=0.21$, $\bar{w}_{\mathrm{linear}}=0.14$.

\paragraph{Interpretation: the saturation form is not robustly
identifiable on the 12-anchor set, but the Hill n=2 form is the
preferred nonlinear form on raw per-step data.}
The apparent contradiction between the synthetic-trace battery
(linear wins) and the raw-trace cross-check (Hill wins) is resolved
by trace length: the 12-anchor traces are 3--30 steps, which is too
short to fit a 2-parameter nonlinear form, so AIC defaults to
linear/null. On 40-step raw per-step traces, the same AIC ranking
favours nonlinear forms, and the sigmoidal Hill n=2 form is
slightly preferred over the exponential (consistent with the
saturating-then-slowing-down shape of an S-curve vs.\ an
asymptote-with-monotonic-curvature exponential). The
implication: the iter17--33 $R_{\max}$ estimates derived under
exponential saturation are not robust to functional-form choice;
the same data is equally or more consistent with a Hill
n=2 form.

\paragraph{Extrapolation comparison (iter37d).}
We re-run the iter21 two-anchor log-log extrapolation battery under
both forms. On the 4-anchor holdout (Qwen3-32B, Qwen3.5-27B,
DeepSeek-V3.1, Qwen3-235B-MoE, Nemotron-120B, Kimi-K2-Thinking),
the mean absolute error on $R_{\max}$ is essentially identical:
$\mathrm{MAE}_{\mathrm{sat}}=0.907$, $\mathrm{MAE}_{\mathrm{hill}}=0.908$.
The $t_{80}$ predictions diverge wildly: $\mathrm{MAE}_{\mathrm{sat}}=584$,
$\mathrm{MAE}_{\mathrm{hill}}=2.4 \times 10^{8}$. The reason is that
on the synthetic frontier traces the saturation fit hits the upper
bound $\lambda \le 5$ and the Hill fit hits the lower bound $K \ge
0.001$, making $t_{80}$ unidentifiable from a 3--30 step trace.
This is itself a \emph{positive} result: the iter17 $t_{80}$
estimates were never identifiable on the frontier set, and
changing the functional form does not change that conclusion.

\paragraph{Recommendation.}
For frontier-trace analysis on $n_{\mathrm{steps}} \le 30$ records,
report the saturation form as the literature default but
\emph{also} report the Hill n=2 fit and a $\Delta\mathrm{AIC}$
column; the two are essentially indistinguishable on this regime.
For $n_{\mathrm{steps}} \ge 30$ per-step data, the Hill n=2 form
is the preferred nonlinear form on the 10 per-step raw traces
measured here, and the exponential saturation is a close second.

\paragraph{Sharpest claim.}
\textit{``The literature's $R(t) = R_{\max}(1 - e^{-\lambda t})$
parameterisation is not the uniquely identified GRPO reward
curve on either the 12-anchor frontier set (where it loses to
linear/null on 9/12 anchors) or the 10-run raw per-step
benchmark (where it loses to the Hill n=2 form on 5/10 runs,
tied on the remaining 5). The iter17--33 R_max estimates should
be reported with a $\pm 0.05$ envelope on the form-choice
sensitivity and a Hill n=2 fit alongside.''}

\paragraph{Artifacts.}
Driver: \texttt{scripts/scaling\_law\_iter37.py},
\texttt{scaling\_law\_iter37b.py},
\texttt{scaling\_law\_iter37c.py},
\texttt{scaling\_law\_iter37d.py}.
Fits: \texttt{scaling\_law\_iter37\_fits.tsv},
\texttt{scaling\_law\_iter37b\_fits.tsv},
\texttt{scaling\_law\_iter37c\_fits.tsv}.
AIC summary: \texttt{scaling\_law\_iter37\_aic.tsv},
\texttt{scaling\_law\_iter37b\_aic.tsv},
\texttt{scaling\_law\_iter37c\_summary.tsv}.
Bootstrap: \texttt{scaling\_law\_iter37\_bootstrap.tsv}.
Extrapolation: \texttt{scaling\_law\_iter37d\_fits.tsv},
\texttt{scaling\_law\_iter37d\_extrap.tsv},
\texttt{scaling\_law\_iter37d\_summary.tsv}.
Figure: \texttt{figures/scaling\_law\_iter37.pdf} (top: stacked
Akaike weights by anchor; bottom: bootstrap win share),
\texttt{figures/scaling\_law\_iter37b.pdf} (Akaike-weight heat-map
on dynamic anchors + bootstrap box-plot),
\texttt{figures/scaling\_law\_iter37c.pdf} (per-run Akaike weights
on raw 40-step per-step traces, GRPO and PPO),
\texttt{figures/scaling\_law\_iter37d.pdf} (extrapolation MAE
saturation vs Hill on 4-anchor holdout).
